\documentclass[11pt]{article}
\usepackage[margin=25mm]{geometry}
\usepackage{fontspec}
\setmainfont{Latin Modern Roman}
\usepackage{amsmath,amssymb,mathtools}
\usepackage{unicode-math}
\setmathfont{Latin Modern Math}
\usepackage{xurl,hyperref}
\hypersetup{colorlinks=true,urlcolor=blue}
\setcounter{secnumdepth}{0}
\setlength{\parindent}{0pt}
\setlength{\parskip}{5pt}
\setlength{\emergencystretch}{3em}
\newcommand{\sref}[2]{\hyperlink{source:#1:#2}{[S#2]}}
\newcommand{\cataloguescope}[1]{\paragraph{Recorded target.}#1}
\newcommand{\R}{\mathbb R}
\newcommand{\C}{\mathbb C}
\newcommand{\N}{\mathbb N}
\newcommand{\Z}{\mathbb Z}
\newcommand{\Q}{\mathbb Q}
\newcommand{\F}{\mathbb F}
\newcommand{\eps}{\varepsilon}
\begin{document}
\section{QQ Systems for Spectral Determinants and Integrable Models}\label{statement}

\subsection{Definitions and mathematical statement}

\paragraph{Shared notation.}$\mathbb N=\{1,2,\ldots\}$; the standard number fields have their usual meanings. All additional conventions are stated locally.


\paragraph{Definitions and conventions.}
An oper is a differential-geometric connection with additional flag and transversality structure; affine opers in the source generalize spectral ordinary differential equations. Spectral determinants encode their spectral or connection data. The $Q\widetilde Q$ system is the source\textquotesingle s functional relation among two families of such functions, indexed by the relevant Lie-algebra data. On the integrable-model side, $Q$-operators or transfer data play an analogous spectral role.

\paragraph{Target retained from the canonical source catalogue.}
\paragraph{Catalogue wording retained for traceability.}
What structural mechanism produces the QQ-tilde system for spectral determinants of affine opers, and is there an analogous system on the integrable-model side?

\paragraph{Source scope and qualification.}
The extraction does not include the functional relation, Lie type, shifts or normalizations. These are source-defined and must be supplied before using this record as a standalone equation; the two structural questions are retained without inventing a new QQ relation.
\subsection{Short English statement}
Explain the origin of the oper spectral-determinant functional relations and find their integrable-model counterpart.
\subsection{Research attempt}
\paragraph{Model, date and investigated scope.}
OpenAI GPT-6 Astra (configured model \texttt{gpt-6-astra}), ultra reasoning; generated 11 October 2026. Proved the rank-two Plucker/Wronskian mechanism as an auxiliary model; neither affine-oper identification nor integrable-model analogue is obtained.

\paragraph{Sources and literature check.}
Read the supplied definition, catalogue wording, source qualifications and reference entries for this question. The historical source/context and inherited status are retained; no fresh literature or novelty certification is claimed. No external theorem is used without proof in this attempt; explicit model assumptions are distinguished from facts about the source model.

\paragraph{A determinant route.} The source qualification does not fix the shifts or Lie type of its \(Q\widetilde Q\) relation. A concrete mechanism worth testing is the determinant identity behind rank-two connection coefficients. Write \([u,v]=u_1v_2-u_2v_1\) for vectors in \(\mathbb C^2\). Direct expansion gives
\[
 [a,b][c,d]-[a,c][b,d]+[a,d][b,c]=0.
\]
For example, when \([a,b]\ne0\), write \(c=\alpha a+\beta b\), \(d=\gamma a+\delta b\). The three terms are respectively \((\alpha\delta-\beta\gamma)[a,b]^2\), \(+\beta\gamma[a,b]^2\), and \(-\alpha\delta[a,b]^2\). They cancel; the same polynomial identity holds without the nonzero assumption by expansion.
\paragraph{Application to a restricted equation.} If \(u,v\) solve \(y''=V(x)y\), their Wronskian \(W(u,v)=uv'-u'v\) satisfies \(W'=uVv-Vuv=0\). Evaluating four solutions at a common point as vectors \((y,y')\) therefore gives the displayed relation with constant Wronskians. It is invariant under changing the evaluation point. However rescaling any solution rescales the corresponding factors: normalizations cannot be ignored when interpreting these as spectral determinants.
\paragraph{First failed identification.} This proves the algebraic identity, not that the source spectral determinants are these Wronskians, that shifted spectral parameters correspond to the four chosen solutions, or that an integrable-model \(Q\)-operator satisfies the same analytic normalization. Those are exactly the missing steps, including higher-rank and affine structure. I retain both source questions and its partially-solved catalogue status; this elementary mechanism supplies no new solution to either.

\paragraph{Verification and adversarial audit.}
Expanded the determinant identity, checked Wronskian constancy only for an equation without a first-derivative term, and retained normalization dependence. This is a GPT-6 Astra Ultra self-audit, not independent human verification. Prompt A drafting followed by prompt B scope/dependency/boundary self-audit. Expanded the determinant identity, checked Wronskian constancy only for an equation without a first-derivative term, and retained normalization dependence. The remaining gap is: Identify normalized source spectral determinants and their parameter shifts with compatible minors, and then construct the integrable-model counterpart. Unresolved status and zero original-target completion are retained.

Coordinator checked all signs in the determinant identity and constancy of Wronskians for the stated second-order equation. Identification with actual spectral determinants remains missing. This additional check was model-assisted, not independent human verification.

\paragraph{Tools, process and contributions.}
Python 3 and local file-reading/writing tools for manuscript serialization and structural checks; mathematical calculations were symbolic reasoning, with no numerical experiment or theorem prover. Contribution prompt A was followed by prompt B. GPT-6 Astra with ultra reasoning drafted the research and performed its stated self-audit; any additional coordinator contribution is identified in the verification and audit. A GPT-6-based Codex coordinator assembled the manuscript and checked repository integration; its exact variant and reasoning configuration were not exposed. No human mathematical contribution, independent human proof check or formal proof certification is supplied. Complete derivations and the consulted source references are included here; private reasoning transcripts are not published.

\paragraph{Final status and first remaining gap.}
\textbf{UNRESOLVED}. Identify normalized source spectral determinants and their parameter shifts with compatible minors, and then construct the integrable-model counterpart.

\section{Completion Estimate}
\noindent\textbf{COMPLETION: 0\%}
The proved calculation is an elementary auxiliary result or a test of a restricted strategy; no new portion of the full source target is established.

\subsection{Sources}
\begin{itemize}
\item[\hypertarget{source:statement:1}{S1}] ULAM.AI and the UnsolvedMath Contributors, UnsolvedMath, catalogue entry. Catalogue curation under CC BY 4.0; definitions and English exposition added locally.
\url{https://huggingface.co/datasets/ulamai/UnsolvedMath}.
\item[\hypertarget{source:statement:2}{S2}] Algebraic Groups (2017), Oberwolfach Reports, DOI 10.4171/owr/2017/21, pp. 1281--1347. Primary report accessed 2026-10-11; relevant PDF pages 50.
\url{https://doi.org/10.4171/owr/2017/21}.
\item[\hypertarget{source:statement:3}{S3}] Frenkel et al., q-Opers and QQ-Systems (catalogue-listed context; not independently consulted in this preparation).
\url{https://arxiv.org/abs/2002.07344}.
\end{itemize}
\label{end:statement}
\end{document}
